\subsection{结果讨论与证据边界}\label{sec:limitations}
\subsubsection{跨任务适用性}
九张历史任务表与本节消融共同显示，接口可移植性和性能泛化应分开理解。SAR-ACD 的正向结果伴随额外 50 轮模块预热，不能隔离纯模块收益；SAMPLE 的 Accuracy 和 Macro-F1 均下降。Ottawa 重跑中 F1 小幅上升，但 Yellow River 四种 batch 的 F1 均下降，且 BS32/64 的 Precision 和 Recall 同时退化。OpenEarthMap-SAR 配置 A 的 mIoU 仅小幅上升，FWIoU/OA 下降，配置 B 三项指标均下降。因此不能将这些实验归纳为所有 SAR 任务的一致改善。

\subsubsection{定性比较与机制验证尚缺}
当前没有纳入绑定同一图像、相同置信度阈值和后处理的经核验预测对比。图~\ref{fig:motivation} 是动机与机制示意，图~\ref{fig:overview} 是计算路径，两者均不是预测可视化。后续应同时展示改善和退化案例，报告固定热图尺度；框内/边缘/背景引导分布、强度匹配和空间干预可用于检验结构依赖，但不能未经实验把学习引导称为可靠性概率或宣称抑制了噪声递推。

\subsubsection{可复现性与下一阶段验收}
随包提供旧 SSDD 原始比例值、新 SADD 回收 JSON、数值计算审计以及完整可编辑表格；原上传底稿逐文件保存在归档目录。完整新 SSDD 控制仍须逐项绑定冻结代码、数据、实际配置、best 权重与独立测试，所有种子和负结果均需报告。未收齐的逐 run 最佳 epoch 不从最大训练预算推测；缺失的结果保持未验收。

文献比较表保留底稿的来源分组和数值，尚未完成主文献逐行数字核对的部分在来源登记中明确列出，不能直接用于跨协议排名。SAR-ShipDataset、SAR-AIRcraft1.0 的公平多种子验证、轻量外部基线和统一效率评估也不由 SADD 12-run 自动替代。本次工作完成的是论文内容整合，不代表这些研究实验全部完成。
